%=======================================================================
%  CONFIGURATION SECURITY AUDIT REPORT — LaTeX TEMPLATE
%=======================================================================
\documentclass[11pt,a4paper]{article}

%------------------------- PACKAGES ------------------------------------
\usepackage[T1]{fontenc}
\usepackage[utf8]{inputenc}
% Falls back to Charter if Latin Modern is not installed
\IfFileExists{lmodern.sty}{\usepackage{lmodern}}%
  {\usepackage{charter}\usepackage[scaled=0.92]{helvet}}
\usepackage[margin=2.4cm,headheight=15pt]{geometry}
\usepackage{xcolor}
\usepackage{graphicx}
\usepackage{booktabs}
\usepackage{array}
\usepackage{longtable}
\usepackage{tabularx}
\usepackage{colortbl}
\usepackage{enumitem}
\usepackage{fancyhdr}
\usepackage{titlesec}
\usepackage{listings}
\usepackage{etoolbox}
\usepackage[most]{tcolorbox}
\usepackage{lastpage}
\usepackage[hidelinks,bookmarks=true]{hyperref}

%------------------------- REPORT METADATA -----------------------------
% >>> EDIT THESE <<<
\newcommand{\ClientName}{ACME Corporation}
\newcommand{\SystemName}{Production Linux Estate \& Perimeter Firewalls}
\newcommand{\EngagementRef}{CSA-2026-014}
\newcommand{\AuditorName}{Jane Doe, Security Engineer}
\newcommand{\AuditorOrg}{Blue Ridge Security Ltd.}
\newcommand{\ReviewerName}{John Smith, Lead Auditor}
\newcommand{\AuditPeriod}{01--12 September 2026}
\newcommand{\ReportDate}{\today}
\newcommand{\DocVersion}{1.0}
\newcommand{\Classification}{CONFIDENTIAL}

%------------------------- SEVERITY COLOURS ----------------------------
\definecolor{sevCritical}{HTML}{7B1113}
\definecolor{sevHigh}{HTML}{C0392B}
\definecolor{sevMedium}{HTML}{D68910}
\definecolor{sevLow}{HTML}{2874A6}
\definecolor{sevInfo}{HTML}{566573}
\definecolor{rulegray}{HTML}{B3B6B7}
\definecolor{lightrow}{HTML}{F2F4F4}
\colorlet{findcolor}{sevInfo}

% Maps a severity name to the colour "findcolor"
\newcommand{\setseverity}[1]{%
  \ifstrequal{#1}{Critical}{\colorlet{findcolor}{sevCritical}}{%
  \ifstrequal{#1}{High}{\colorlet{findcolor}{sevHigh}}{%
  \ifstrequal{#1}{Medium}{\colorlet{findcolor}{sevMedium}}{%
  \ifstrequal{#1}{Low}{\colorlet{findcolor}{sevLow}}{%
  \colorlet{findcolor}{sevInfo}}}}}}

% Inline coloured severity badge: \sev{High}
\newcommand{\sev}[1]{\setseverity{#1}%
  \textcolor{white}{\colorbox{findcolor}{\strut\footnotesize\bfseries\,#1\,}}}

%------------------------- FINDING ENVIRONMENT -------------------------
% Usage: \begin{finding}{FIND-001}{Short title}{High} ... \end{finding}
\newenvironment{finding}[3]{%
  \setseverity{#3}%
  \phantomsection\addcontentsline{toc}{subsection}{#1 --- #2}%
  \begin{tcolorbox}[
    breakable, enhanced, sharp corners=downhill,
    colback=white, colframe=findcolor, boxrule=0.8pt,
    left=8pt, right=8pt, top=6pt, bottom=6pt,
    fonttitle=\bfseries\color{white},
    title={#1 \; --- \; #2 \hfill \MakeUppercase{#3}}]%
}{\end{tcolorbox}\vspace{4pt}}

% Metadata block used at the top of each finding
\newcommand{\findingmeta}[4]{%
  \begin{tabularx}{\linewidth}{@{}>{\bfseries}p{3.1cm}X@{}}
    Affected assets & #1 \\
    Control mapping & #2 \\
    Likelihood / Impact & #3 \\
    Status & #4 \\
  \end{tabularx}\vspace{4pt}\hrule\vspace{6pt}}

\newcommand{\fsec}[1]{\vspace{4pt}\noindent\textbf{\textcolor{findcolor}{#1}}\par\nobreak\vspace{2pt}}

%------------------------- CODE / CONFIG LISTINGS ----------------------
\definecolor{lstbg}{HTML}{F7F8F9}
\lstdefinestyle{config}{
  backgroundcolor=\color{lstbg}, basicstyle=\ttfamily\footnotesize,
  breaklines=true, frame=leftline, framerule=1.2pt, rulecolor=\color{rulegray},
  xleftmargin=8pt, framexleftmargin=8pt, showstringspaces=false,
  commentstyle=\color{sevInfo}\itshape, keywordstyle=\color{sevLow}\bfseries,
  captionpos=b, aboveskip=6pt, belowskip=6pt}
\lstset{style=config}

%------------------------- HEADERS / FOOTERS ---------------------------
\pagestyle{fancy}\fancyhf{}
\renewcommand{\headrulewidth}{0.4pt}
\renewcommand{\footrulewidth}{0.4pt}
\fancyhead[L]{\footnotesize\ClientName\ --- Configuration Security Audit}
\fancyhead[R]{\footnotesize\textbf{\Classification}}
\fancyfoot[L]{\footnotesize Ref: \EngagementRef\ \textbar\ v\DocVersion}
\fancyfoot[R]{\footnotesize Page \thepage\ of \pageref{LastPage}}

\titleformat{\section}{\Large\bfseries}{\thesection}{0.7em}{}[\vspace{2pt}\hrule]
\setlist[itemize]{leftmargin=1.4em,itemsep=2pt,topsep=3pt}
\setlist[enumerate]{leftmargin=1.6em,itemsep=2pt,topsep=3pt}
\renewcommand{\arraystretch}{1.25}

%=======================================================================
\begin{document}
%------------------------- TITLE PAGE ----------------------------------
\begin{titlepage}
\centering
\vspace*{1.5cm}
{\small\bfseries\Classification}\\[2cm]
% \includegraphics[width=5cm]{logo.png}\\[1cm]
{\Huge\bfseries Configuration Security Audit\par}
\vspace{0.6cm}
{\LARGE \SystemName\par}
\vspace{2cm}
\begin{tabular}{@{}>{\bfseries}l l@{}}
\toprule
Client & \ClientName \\
Engagement reference & \EngagementRef \\
Audit period & \AuditPeriod \\
Prepared by & \AuditorName \\
Organisation & \AuditorOrg \\
Reviewed by & \ReviewerName \\
Report date & \ReportDate \\
Version & \DocVersion \\
Classification & \Classification \\
\bottomrule
\end{tabular}
\vfill
{\footnotesize This document contains confidential information about the security posture of
\ClientName. Distribution is restricted to the parties listed in the distribution table.
Do not forward, copy, or store outside approved systems.\par}
\end{titlepage}

%------------------------- DOCUMENT CONTROL ----------------------------
\section*{Document Control}
\addcontentsline{toc}{section}{Document Control}

\noindent\textbf{Revision history}
\begin{longtable}{@{}l l p{5.6cm} l l@{}}
\toprule
\textbf{Version} & \textbf{Date} & \textbf{Summary of change} & \textbf{Author} & \textbf{Approved} \\
\midrule\endhead
0.1 & 2026-09-08 & Initial draft & J. Doe & --- \\
0.9 & 2026-09-11 & Peer review comments applied & J. Doe & J. Smith \\
1.0 & 2026-09-12 & Final issue to client & J. Doe & J. Smith \\
\bottomrule
\end{longtable}

\noindent\textbf{Distribution list}
\begin{longtable}{@{}l l l@{}}
\toprule
\textbf{Name} & \textbf{Role} & \textbf{Organisation} \\
\midrule\endhead
A. Example & CISO & \ClientName \\
B. Example & Head of Infrastructure & \ClientName \\
\ReviewerName & Lead Auditor & \AuditorOrg \\
\bottomrule
\end{longtable}

\newpage
\tableofcontents
\newpage

%=======================================================================
\section{Executive Summary}
Summarise, in language a non-technical reader can act on: what was assessed,
how it was assessed, the overall state of the configuration baseline, and the
two or three issues that matter most. Keep this section to one page.

\vspace{6pt}
\noindent\textbf{Overall assessment:} \emph{e.g.\ Moderate risk --- the estate is
broadly hardened, but inconsistent baseline enforcement on newly provisioned
hosts leaves a recurring gap.}

\vspace{10pt}
\noindent\textbf{Findings by severity}
\begin{center}
\begin{tabular}{@{}l c c c c c@{}}
\toprule
& \sev{Critical} & \sev{High} & \sev{Medium} & \sev{Low} & \sev{Info} \\
\midrule
Count & 1 & 3 & 5 & 4 & 2 \\
\bottomrule
\end{tabular}
\end{center}

\vspace{6pt}
\noindent\textbf{Key themes}
\begin{enumerate}
  \item \textbf{Theme one.} One or two sentences on the systemic cause, not just the symptom.
  \item \textbf{Theme two.} Where the same misconfiguration recurs across assets, say so here.
  \item \textbf{Theme three.} Note anything done well --- it tells the reader what to protect.
\end{enumerate}

%=======================================================================
\section{Scope and Objectives}

\subsection{Objectives}
\begin{itemize}
  \item Verify that deployed configurations match the approved hardening baseline.
  \item Identify misconfigurations that weaken authentication, authorisation, logging, or network exposure.
  \item Assess configuration change control and drift detection.
  \item Map gaps to the applicable control framework and provide prioritised remediation.
\end{itemize}

\subsection{Assets in scope}
\begin{longtable}{@{}l l l l l@{}}
\toprule
\textbf{Asset ID} & \textbf{Hostname / Identifier} & \textbf{Type} & \textbf{Version} & \textbf{Environment} \\
\midrule\endhead
AST-001 & web-prod-01 & Ubuntu Server & 24.04 LTS & Production \\
AST-002 & fw-edge-01 & Firewall & PAN-OS 11.1 & Production \\
AST-003 & db-prod-01 & PostgreSQL & 16.3 & Production \\
AST-004 & k8s-prod & Kubernetes cluster & 1.30 & Production \\
\bottomrule
\end{longtable}

\subsection{Out of scope}
\begin{itemize}
  \item Application source code review and dynamic application testing.
  \item Physical security and personnel controls.
  \item Third-party SaaS tenants not administered by \ClientName.
\end{itemize}

\subsection{Assumptions and limitations}
\begin{itemize}
  \item Configuration was reviewed as at \AuditPeriod; later changes are not reflected.
  \item Read-only credentials were provided; no exploitation was attempted.
  \item Findings are based on the evidence supplied; undisclosed compensating controls may alter risk ratings.
\end{itemize}

%=======================================================================
\section{Methodology}
Describe how the review was performed so a third party could repeat it.

\subsection{Approach}
\begin{enumerate}
  \item \textbf{Collection} --- configuration exports, running state, and policy documents gathered from each in-scope asset.
  \item \textbf{Baseline comparison} --- settings compared against the benchmarks listed below and against \ClientName's internal standard.
  \item \textbf{Manual review} --- automated results triaged to remove false positives and to assess contextual risk.
  \item \textbf{Validation} --- deviations confirmed with the system owner before rating.
  \item \textbf{Rating and reporting} --- each deviation rated on likelihood and impact, then mapped to controls.
\end{enumerate}

\subsection{Reference standards}
\begin{longtable}{@{}l p{7.5cm} l@{}}
\toprule
\textbf{Ref} & \textbf{Standard / benchmark} & \textbf{Version} \\
\midrule\endhead
CIS & CIS Benchmark for the relevant platform & v2.0.0 \\
NIST & NIST SP 800-53 Rev.\ 5 security controls & Rev.\ 5 \\
ISO & ISO/IEC 27001:2022 Annex A & 2022 \\
INT & \ClientName\ Internal Hardening Standard & v3.1 \\
\bottomrule
\end{longtable}

\subsection{Tooling}
\begin{longtable}{@{}l p{6.5cm} l@{}}
\toprule
\textbf{Tool} & \textbf{Purpose} & \textbf{Version} \\
\midrule\endhead
Tool A & Automated benchmark scanning & 1.2.3 \\
Tool B & Configuration export and diff & 4.5 \\
Manual & Shell inspection and interview & --- \\
\bottomrule
\end{longtable}

\subsection{Severity rating}
Severity is derived from likelihood and impact as set out in Appendix~\ref{app:ratings}.

%=======================================================================
\section{Summary of Findings}
\begin{longtable}{@{}l p{6.2cm} l l l@{}}
\toprule
\textbf{ID} & \textbf{Finding} & \textbf{Severity} & \textbf{Assets} & \textbf{Status} \\
\midrule\endhead
\rowcolor{lightrow}
CSA-001 & Password authentication permitted for remote administrative access & \sev{Critical} & AST-001 & Open \\
CSA-002 & Permissive inbound rule exposes management interface & \sev{High} & AST-002 & Open \\
CSA-003 & Audit logging disabled for privileged operations & \sev{Medium} & AST-003 & Open \\
CSA-004 & Legacy protocol version still enabled & \sev{Low} & AST-004 & Accepted \\
\bottomrule
\end{longtable}

%=======================================================================
\section{Detailed Findings}

% ----------------------- TEMPLATE FINDING ----------------------------
\begin{finding}{CSA-001}{Password authentication permitted for remote administrative access}{Critical}
\findingmeta
  {AST-001 (web-prod-01)}                    % affected assets
  {CIS 5.2.x; NIST IA-2, AC-17; ISO A.8.5}   % control mapping
  {Likely / Severe}                          % likelihood / impact
  {Open}                                     % status

\fsec{Description}
State plainly what the configuration is and how it deviates from the baseline.
Avoid tool output pasted without interpretation.

\fsec{Evidence}
\begin{lstlisting}[caption={Excerpt from /etc/ssh/sshd\_config on web-prod-01},label={lst:csa001}]
# Observed setting
PasswordAuthentication yes
PermitRootLogin yes
# Expected per baseline INT-3.1 section 4.2
# PasswordAuthentication no
# PermitRootLogin prohibit-password
\end{lstlisting}

\fsec{Risk and impact}
Explain the realistic consequence in business terms: what an attacker gains,
which data or service is reachable, and what else breaks as a result.

\fsec{Recommendation}
\begin{enumerate}
  \item Set the control to its required value, quoting the exact directive or console path.
  \item Enforce the setting through configuration management so it survives rebuilds.
  \item Add a detection rule or scheduled check for drift on this setting.
\end{enumerate}

\fsec{Remediation effort}
Low --- single configuration change plus service reload; no downtime expected.

\fsec{References}
CIS Benchmark section x.y.z; vendor hardening guide; internal standard INT-3.1 \S4.2.
\end{finding}

% ----------------------- COPY FROM HERE FOR NEW FINDINGS -------------
\begin{finding}{CSA-002}{Short, factual finding title}{High}
\findingmeta{AST-00x}{Framework control IDs}{Possible / Major}{Open}

\fsec{Description}
Replace with the observed configuration and the deviation.

\fsec{Evidence}
\begin{lstlisting}[caption={Source of evidence --- file, command, or console screen}]
key = observed_value    # expected: required_value
\end{lstlisting}

\fsec{Risk and impact}
Replace.

\fsec{Recommendation}
\begin{enumerate}
  \item Replace.
\end{enumerate}

\fsec{Remediation effort}
Replace.

\fsec{References}
Replace.
\end{finding}

%=======================================================================
\section{Remediation Roadmap}
Prioritise by risk reduction per unit of effort, not by severity alone.

\begin{longtable}{@{}l p{5.4cm} l l l@{}}
\toprule
\textbf{ID} & \textbf{Action} & \textbf{Priority} & \textbf{Owner} & \textbf{Target date} \\
\midrule\endhead
CSA-001 & Disable password authentication and enforce via config management & Immediate ($\leq$7 days) & Infrastructure & 2026-09-21 \\
CSA-002 & Restrict management interface to jump-host source range & Short term ($\leq$30 days) & Network & 2026-10-12 \\
CSA-003 & Enable privileged operation auditing and forward to SIEM & Medium term ($\leq$90 days) & Platform & 2026-12-10 \\
\bottomrule
\end{longtable}

\subsection{Systemic recommendations}
\begin{itemize}
  \item Bring all in-scope hosts under a single enforced configuration baseline.
  \item Run benchmark scanning on a schedule and alert on drift rather than auditing point-in-time.
  \item Require baseline compliance as a gate in the provisioning pipeline.
\end{itemize}

%=======================================================================
\section{Conclusion}
Close with the overall posture, the residual risk once the roadmap is complete,
and any recommended date for re-testing or follow-up review.

%=======================================================================
\appendix
\section{Severity and Rating Definitions}\label{app:ratings}

\begin{longtable}{@{}l p{9.5cm}@{}}
\toprule
\textbf{Severity} & \textbf{Definition} \\
\midrule\endhead
\sev{Critical} & Directly exploitable, exposes sensitive data or privileged access; remediate immediately. \\
\sev{High} & Significant weakening of a security control; exploitable with modest effort or conditions. \\
\sev{Medium} & Meaningful gap requiring chained conditions or limited to a narrower asset set. \\
\sev{Low} & Minor deviation with limited practical impact; hardening improvement. \\
\sev{Info} & No direct risk; observation, good practice, or note for future work. \\
\bottomrule
\end{longtable}

\vspace{6pt}
\noindent\textbf{Likelihood / impact matrix}
\begin{center}
\begin{tabular}{@{}l c c c@{}}
\toprule
\textbf{Likelihood $\downarrow$ / Impact $\rightarrow$} & \textbf{Minor} & \textbf{Major} & \textbf{Severe} \\
\midrule
Likely   & \sev{Medium} & \sev{High}   & \sev{Critical} \\
Possible & \sev{Low}    & \sev{Medium} & \sev{High} \\
Unlikely & \sev{Info}   & \sev{Low}    & \sev{Medium} \\
\bottomrule
\end{tabular}
\end{center}

%=======================================================================
\section{Control Framework Mapping}
\begin{longtable}{@{}l p{4.6cm} l l l@{}}
\toprule
\textbf{Finding} & \textbf{Control area} & \textbf{CIS} & \textbf{NIST 800-53} & \textbf{ISO 27001} \\
\midrule\endhead
CSA-001 & Remote access authentication & 5.2.x & IA-2, AC-17 & A.8.5 \\
CSA-002 & Network boundary protection & 3.x & SC-7 & A.8.20 \\
CSA-003 & Audit and accountability & 4.x & AU-2, AU-12 & A.8.15 \\
\bottomrule
\end{longtable}

%=======================================================================
\section{Evidence Index}
\begin{longtable}{@{}l p{5.6cm} l l@{}}
\toprule
\textbf{Ref} & \textbf{Artefact} & \textbf{Source} & \textbf{Collected} \\
\midrule\endhead
EV-001 & sshd\_config export & web-prod-01 & 2026-09-03 \\
EV-002 & Firewall policy export & fw-edge-01 & 2026-09-04 \\
EV-003 & Benchmark scan report & Tool A & 2026-09-05 \\
\bottomrule
\end{longtable}

\noindent\footnotesize Evidence is retained in accordance with the engagement
agreement and destroyed on the agreed schedule.

\end{document}